\clearpage % Start a new page
\setstretch{1.5} % Set the line spacing to 1.5, this makes the following tables easier to read
\chapter*{Lista de símbolos}
\addcontentsline{toc}{chapter}{Lista de Símbolos}
\lhead{\emph{Lista de símbolos}}
    \begin{longtable}{lp{0.84\linewidth}}
    
    %	Darle mejor ancho a la segunda columna 
        $\epsilon$ & Una serie temporal especial conocida como
ruido blanco (\textit{white noise}, en inglés). \\

        $\phi_i$ & Representa a variables que deben resolverse  en los modelos de autoregresión. \\

        $\delta$ & Es una constante empleada en los modelos de autoregresión o ARIMA. \\

    $y_i$ & Representa a cada uno de los valores de la serie temporal. \\

    $\mu$ & Es una constante empleada en los modelos de media móvil. \\

    $\theta_i$ & Representa a variables que deben resolverse  en los modelos de media móvil o ARIMA. \\

     $c$ & Es una constante empleada en los modelos de ARIMA. \\

    $i_{t}$ & Representa a la puerta de entrada en modelos LSTM. \\

    $o_{t}$ & Representa a la puerta de salida en modelos LSTM. \\

    $f_{t}$ & Representa a la puerta de olvido en modelos LSTM. \\

    $c_{t-new}$ & Es el nuevo estado, candidato, de la célula en modelos LSTM. \\

    $c_{t}$ & Es vector que representan el estado final de la célula en modelos LSTM. \\

    $h_{t}$ & Es el vector de la capa oculta que sale de una unidad recurrente. \\

    $W, U, b$ & Son matrices cuyos valores deben ser aprendidos durante el proceso de entrenamiento de las redes recurrentes. \\

    $\odot$ & Representa a la operación  denominada \textit{producto de Hadamard}. \\

    $logistic$ & Representa a la función de activación \textit{logistic}. \\

    $tanh$ & Representa a la función de activación \textit{tangente hiperpólica}. \\

    $z_{t}$ & Representa a la puerta de actualización en modelos GRU. \\

    $r_{t}$ & Representa a la puerta de reseteo en modelos GRU. \\

    $h^-_{t}$ & Es un vector cuyo valor es utilizado para calcular \(h_{t}\) en el modelado con GRU. \\

    $x_i'$ & Representa al resultado de aplicar la función de normalización \textit{min-max}. \\

    $o_i$ & Representa a un valor del conjunto de valores observados. \\

    $o_{min}$ & Representa el valor mínimo del conjunto de valores observados. \\

    $o_{max}$ & Representa el valor máximo del conjunto de valores observados. \\

    $\overline{o}$ & Representa el valor promedio del conjunto de valores observados. \\

    $p_i$ & Representa a un valor del conjunto de valores pronosticados. \\

    $\overline{p}$ & Representa el valor promedio del conjunto de valores pronosticados. \\

    $N$ & Representa a la cantidad total de elementos en un conjunto de datos.\\

        \\\hline 
    \end{longtable}